% SOURCE_AUTHORITY: sga5_fr_workpass.tex, lines 15248--15264 (LF-normalized unit).
% SOURCE_SHA256: 791F4EFFC5E02832D5D77ED03518C8156D6F07E4C8238B03545DB93D883FBB28
\begin{statement}{Corolário}
Suponhamos que o ideal $I$ seja nilpotente. Para todo complexo acíclico $N_0$ de $A_0$-módulos projetivos, limitado à direita, existe um complexo acíclico $N$ de $A$-módulos projetivos, limitado à direita, tal que $T(N)\simeq N_0$.
\end{statement}

Como o complexo $N_0$ é acíclico, decompõe-se em sequências exatas curtas
\[
0\longrightarrow Z_0^i\longrightarrow N_0^i\longrightarrow Z_0^{i+1}\longrightarrow0,
\]
onde $Z_0^i$ designa, para cada $i$, o submódulo dos ciclos de $N_0^i$; além disso, $N_0^i=0$ para $i$ suficientemente grande. Por indução descendente sobre $i$, vê-se que as sequências exatas precedentes são cindidas e que $Z_0^i$ é projetivo para todo $i$. Levanta-se $N_0$ levantando sucessivamente todas essas sequências exatas, por indução descendente sobre $i$. Trata-se, portanto, de demonstrar que, se
\[
S_0:0\longrightarrow P_0\longrightarrow Q_0\longrightarrow R_0\longrightarrow0
\]
é uma sequência exata de $A_0$-módulos, com $Q_0$ projetivo, e se $R$ é um $A$-módulo tal que $T(R)\simeq R_0$, existe uma sequência exata
\[
S:0\longrightarrow P\longrightarrow Q\longrightarrow R\longrightarrow0
\]
de $A$-módulos, com $Q$ projetivo, tal que $T(S)\simeq S_0$. Ora, pode-se levantar $Q_0$ a um $A$-módulo projetivo $Q$ pelo lema 3(ii), e depois $Q_0\to R_0$ a uma aplicação $A$-linear $Q\to R$ pelo lema 3(i). Esta aplicação é sobrejetiva pelo lema de Nakayama; põe-se $P=\ker(Q\to R)$ para completar a sequência exata $S$.
